\section{The actual original General history crossing an interior empty step}

The original general invariant keeps its actual nonmonomial order and
the companion formed at the start of its original order regime.  To
compare an interior empty insertion, that order and its whole selected
history must be computed from the original datum, including an empty
tail.  The following are genuine history and canonical-input producers
for the positive General branch.  They do not assume a companion value
or new canonical-history comparison.  The source route is W\l odarczyk's
General step and Bierstone--Milman's trivial-blowup convention
\cite{wlodarczyk05-hironaka,bierstone-milman08-functoriality}.

\begin{theorem}[Original nonmonomial order on every interior empty history]
\label{hir:original-interior-empty-all-tail-order}
\source{wlodarczyk05-hironaka,bierstone-milman08-functoriality}
\lean{Hironaka.HironakaOriginalGeneralCrossingHistory.nmOrd_interior_empty_all_tails}
\leanok
\uses{hir:original-interior-empty-history,hir:original-empty-boundary-test-history,bcalc:nonmonomial-order-transport}
Let $M$ be originally well formed, $H$ an actual test history of $M$,
$S$ any actual subsequent history, and $y$ any actual final point.  Then
the nonmonomial orders of the SAME original $M$ on $H\epsilon S$ and $HS$
at $y$ agree.  No nonempty-tail, positive-order or desired equality
condition is required.
\end{theorem}
\begin{proof}
The actual transformed head $K=M_H$ is well formed.  Its genuine identity
empty blowup is admissible, hence has a coherent transformed ideal.
Removing its new empty boundary entry leaves the actual empty-step base
$K_\epsilon^{\mathrm{base}}$, coherently restricted-equivalent to $K$
by the genuine identity pullback dictionary.  Their carriers are the same
closed original carrier, and both ideals are coherent.  Restricted
equivalence therefore gives equal nonmonomial orders after the entire
actual suffix $S$.  The actual new empty boundary entry changes no
nonmonomial order, also after $S$; its index is internally valid because
the base boundary has exactly $|K.E|$ members.  The real transform-append
and transform-step dictionaries identify these two computed transforms
with $M_{H\epsilon S}$ and $M_{HS}$.  This proves the stated original
order equality for every $S$, including the immediate empty step.
\end{proof}

\begin{theorem}[The actual original order-regime run across insertion]
\label{hir:original-interior-empty-order-regime-run}
\source{wlodarczyk05-hironaka,bierstone-milman08-functoriality}
\lean{Hironaka.HironakaOriginalGeneralCrossingHistory.ordRegime_dictionary,Hironaka.HironakaOriginalGeneralCrossingHistory.ordRegime_run_dictionary}
\leanok
\uses{hir:original-interior-empty-all-tail-order,hir:tower-split}
For these actual data and any fixed original order $o$, put
$r=\operatorname{run}_{M,o}(HS,y)$.  The actual run on $H\epsilon S$
is $r$ if $r<|S|$, is $r+1$ if $r\geq|S|$ and the original point
$S.\operatorname{proj}(y)$ is in order regime $o$ at $H$, and is $|S|$
otherwise.  The regime is the literal original transformed-cosupport
membership together with $\nu=o$.
\end{theorem}
\begin{proof}
Induct on the actual last steps of $S$.  At the empty tail, the run
recursion on the genuine identity empty step reads the original regime
at $H$: it adds one to the original run if that regime holds, and is
zero otherwise.  This is exactly the displayed formula when $|S|=0$.
At a nonempty last step, the preceding theorem identifies the original
order regime at the two predecessor points.  If the predecessor fails
the regime, both runs are zero and the old run is shorter than the
positive tail length.  If it satisfies the regime, each actual run is
the predecessor run plus one.  Apply induction and compare the old
predecessor run with the predecessor tail length.  Each of the three
inductive cases gives precisely the corresponding stated current case.
The actual whole cosupport dictionary supplies the cosupport component
of the two predecessor regimes; the order dictionary supplies their
order component.  The composite projections identify the original base
point throughout.
\end{proof}

\begin{theorem}[Both whole General splits retain the original assignment head]
\label{hir:original-interior-general-crossing-splits}
\source{wlodarczyk05-hironaka,bierstone-milman08-functoriality}
\lean{Hironaka.HironakaOriginalGeneralCrossingHistory.actual_general_crossing_splits}
\leanok
\uses{hir:original-interior-empty-order-regime-run,hir:original-crossing-run-base-regime,hir:original-crossing-restricted-empty-history}
Assume $y$ is in the actual original transformed cosupport, its actual
current order is $o$, and the original run $r$ reaches
$S$.  The actual new order-regime run is $r+1$.  Put $a=r-|S|$,
$D=\operatorname{stage}(H,a)$, $K=\operatorname{head}(H,a)$ and
$B=\operatorname{tail}(H,a)$.  Both entire selected General split records
have the SAME original stage/head $(D,K)$, with tails $BS$ and
$B\epsilon S$.  Hence the actual original order mask
$U=\operatorname{ordMask}(M_K,o)$ and WHOLE companion $J$ formed from
$M,o,K$ are identical at the two assignments.  The new restricted
companion tail has the genuine restricted-interior-empty diagram.
\end{theorem}
\begin{proof}
Actual cosupport membership and current order supply current order-regime membership.  A
run reaching $S$ internally supplies membership at the original $H$
point by the true run-reaching induction.  The preceding run formula
therefore gives $r+1$.  The actual full split formula above the tail
length gives the old split $(D,K,BS)$.  Apply it to $H\epsilon,S$ at
$r+1$: the remaining split count is $a+1$, whose first removed step is
the identity empty blowup.  The actual split recursion therefore gives
$(D,K,B\epsilon S)$.  Both lower constructions use literally the same
$M,o,K$, so their order mask and entire companion agree without a chosen
comparison hypothesis.  Restrict the real empty-history diagram to this
literal original open $U$; its identity base and actual final open
identifications give the entire asserted tail isomorphism.
\end{proof}

\begin{theorem}[Original full canonical input supplies the crossing companion data]
\label{hir:original-canonical-general-crossing-companion-data}
\source{wlodarczyk05-hironaka,bierstone-milman08-functoriality}
\lean{Hironaka.HironakaOriginalGeneralCrossingHistory.original_canonical_general_crossing_data,Hironaka.HironakaOriginalGeneralCrossingHistory.original_canonical_general_companion_history}
\leanok
\uses{hir:original-interior-general-crossing-splits,hir:general-canonical-tail,hir:original-interior-empty-test-history}
Let $n>0$, let $M$ belong to the original dimension-$(k+1)$ invariant
domain, and let $HS$ be canonical for its actual full $R_{k+1}$.
Let $y$ be an original transformed-cosupport point with actual positive
order $o$ and actual order-regime run reaching $S$.  Use the true original
$K,B,U,J$ just computed.  Original data internally supply mask membership,
the true dimension-$(k+1)$ domain and maximal order of $J$, the entire
OLD restricted companion tail canonical for the actual Step One value
formed from $R_k$, and its actual cosupport point.  They also supply the
WHOLE new restricted companion tail diagram and its genuine testness.
No companion domain, maximal-order, test, canonical-history, mask or value
comparison is a caller premise.
\end{theorem}
\begin{proof}
The actual original family supplies testness of its canonical $HS$ and
hence of $H$.  Its positive branch is the genuine General value whose
inner maximal-order value is the actual Step One construction from
$R_k$.  That construction has its proved actual maximal-order value
properties.  The General canonical-tail theorem therefore gives actual
companion lower data at the original order-regime split.  Substitute
the preceding WHOLE computed original split before unpacking the
dependent mask member and cosupport point.  This gives exactly the
domain, maximal order, canonical restricted old tail and cosupport
stated for the literal $U,J,BS$.  The preceding new split retains that
same original $U,J$ and gives its genuine whole restricted empty-history
diagram.  The true old canonical companion tail is a test history.
Normalize its restriction of concatenation by the actual tower
isomorphism; testness transfers and the real identity-pullback gives
the same coherent $J$.  The original interior-empty test equivalence
supplies testness of the normalized inserted companion tail.  The true
new restriction diagram transfers this testness back to the whole
new restricted companion tail.  These constructions use only the
internally supplied original companion data and actual geometry.
\end{proof}

The added interior companion history is not asserted to be canonical or
to have the required positive recursive value yet.  Full positive inner
insertion, outer General values, mixed spectra and noncompact scheduling
remain further obligations.
